\documentclass[UTF8,zihao=-4]{ctexart}
\usepackage[a4paper,margin=1.65cm]{geometry}
\usepackage{amsmath,amssymb,bm}
\usepackage{booktabs,longtable,array}
\usepackage{xcolor}
\usepackage{hyperref}
\usepackage{enumitem}

\hypersetup{colorlinks=true,linkcolor=blue!60!black,urlcolor=blue!60!black}
\setlength{\parindent}{0pt}
\setlength{\parskip}{0.45em}
\renewcommand{\arraystretch}{1.18}
\setlist[enumerate]{itemsep=0.25em,topsep=0.25em}

\newcommand{\dd}{\,\mathrm d}
\newcommand{\ii}{\mathrm i}
\newcommand{\e}{\mathrm e}
\newcommand{\key}[1]{\textcolor{blue!65!black}{#1}}
\newcommand{\note}[1]{\textcolor{red!65!black}{#1}}

\title{\bfseries 第5章理想流体：原题练习汇总完整解答版}
\author{}
\date{}

\begin{document}
\maketitle

本文件对应《第5章理想流体：原题练习汇总（去作业版）》。
写法按“考试可抄写”的顺序整理：先列方程或基本解，再代入题目条件，最后给结论。
符号约定：平面流动复势 \(w=\varphi+\ii\psi\)，点源强度 \(Q>0\) 表示流出，点涡强度用 \(\Gamma\) 或 \(q\) 表示，空气/液体密度按题给。

\section*{一、必用公式}

\begin{longtable}{p{3.2cm}p{11.5cm}}
\toprule
项目 & 公式 \\
\midrule
平面不可压流函数 & \(u=\dfrac{\partial \psi}{\partial y},\quad v=-\dfrac{\partial \psi}{\partial x}\)，流线为 \(\psi=C\)。\\
速度势 & \(u=\dfrac{\partial \varphi}{\partial x},\quad v=\dfrac{\partial \varphi}{\partial y}\)，无旋时 \(\nabla^2\varphi=0\)。\\
不可压条件 & \(\dfrac{\partial u}{\partial x}+\dfrac{\partial v}{\partial y}=0\)。\\
无旋条件 & \(\omega_z=\dfrac{\partial v}{\partial x}-\dfrac{\partial u}{\partial y}=0\)。\\
平面点源 & \(\displaystyle w=\frac{Q}{2\pi}\ln(z-z_0)\)。\\
平面点涡 & \(\displaystyle w=-\frac{\ii\Gamma}{2\pi}\ln(z-z_0)\)。\\
均匀流 & 沿 \(+x\)：\(w=Uz\)，沿方向角 \(\alpha\)：\(w=U\e^{-\ii\alpha}z\)。\\
圆柱绕流 & \(\displaystyle w=U\left(z+\frac{a^2}{z}\right)-\frac{\ii\Gamma}{2\pi}\ln z\)。\\
物面条件 & 不可穿透：物面是流线，\(\psi=\text{常数}\)。\\
\bottomrule
\end{longtable}

\section*{二、课本第5章习题完整解答}

\subsection*{5.1 速度势建立无界流体中椭圆柱运动问题}

椭圆边界为
\[
F(x,y)=\frac{x^2}{a^2}+\frac{y^2}{b^2}-1=0 .
\]
流体理想、不可压、无旋，所以速度势满足拉普拉斯方程：
\[
\nabla^2\varphi=0,\qquad (x,y)\ \text{在椭圆外部}.
\]

物面速度由平动和转动叠加：
\[
\bm V_b=U_0(t)\bm i+\omega(t)\bm k\times (x\bm i+y\bm j)
=\bigl(U_0(t)-\omega y\bigr)\bm i+\omega x\bm j .
\]
椭圆外法向量方向可取
\[
\bm n \parallel \nabla F=\left(\frac{2x}{a^2},\frac{2y}{b^2}\right).
\]
不可穿透条件为
\[
\frac{\partial \varphi}{\partial n}=\bm V_b\cdot \bm n
\quad \text{在}\quad \frac{x^2}{a^2}+\frac{y^2}{b^2}=1.
\]
无界远处流体原静止，故
\[
\nabla\varphi\to 0,\qquad r\to\infty .
\]
题设环量 \(\Gamma=0\)，补充条件为
\[
\oint_C \nabla\varphi\cdot \dd\bm l=0 .
\]
这四条就是该问题的完整数学提法。

\subsection*{5.2 用流函数建立管内障碍物绕流封闭方程组}

不可压平面无旋流动，流函数满足
\[
\nabla^2\psi=0
\]
在流体区域内成立。

管壁和固体边界均不可穿透，所以它们都是流线：
\[
\psi=C_1\quad \text{上壁},\qquad
\psi=C_2\quad \text{下壁},\qquad
\psi=C_3\quad \text{固体表面}.
\]
若管内从左向右的体积流量为 \(Q\)，则上、下壁两条流线的差为
\[
C_1-C_2=Q .
\]
入口截面速度均匀且沿 \(x\) 轴，设入口高度为 \(2h\)，则
\[
u=\frac{Q}{2h},\qquad v=0 .
\]
由 \(u=\partial\psi/\partial y\) 得入口处
\[
\psi(y)=\frac{Q}{2h}y+C .
\]
出口处只需给出与入口相同的总流量条件。以上方程、壁面条件和流量条件构成封闭方程组。

\subsection*{5.3 用流函数建立运动椭圆柱绕流问题}

坐标取在物体上。流体绝对运动无旋，因此在随体坐标中仍可用流函数表示相对流动：
\[
\nabla^2\psi=0 .
\]
物面不可穿透，椭圆面为流线：
\[
\psi=C_b,\qquad \frac{x^2}{a^2}+\frac{y^2}{b^2}=1 .
\]
远处流体原静止，而坐标系随物体以 \(U_0(t)\) 平移，所以在随体坐标中远处来流为
\[
u\to -U_0(t),\qquad v\to 0,\qquad r\to\infty .
\]
对应流函数远处渐近为
\[
\psi\sim -U_0(t)y,\qquad r\to\infty .
\]
若已知环量为 \(\Gamma\)，还需附加
\[
\oint_C \bm V\cdot \dd\bm l=\Gamma .
\]
这就是用流函数建立的定解问题。

\subsection*{5.4 判断流函数、速度势、变形率和刚体转动角速度}

题给
\[
u=x^2y+y^2,\qquad v=x^2-xy^2,\qquad w=0 .
\]

\textbf{1. 判断是否存在不可压流函数}

\[
\frac{\partial u}{\partial x}=2xy,\qquad
\frac{\partial v}{\partial y}=-2xy .
\]
所以
\[
\frac{\partial u}{\partial x}+\frac{\partial v}{\partial y}=2xy-2xy=0.
\]
连续方程满足，因此存在平面不可压流函数 \(\psi\)。

\textbf{2. 求流函数}

由
\[
u=\frac{\partial\psi}{\partial y}=x^2y+y^2
\]
对 \(y\) 积分：
\[
\psi=\frac12x^2y^2+\frac13y^3+f(x).
\]
再由
\[
v=-\frac{\partial\psi}{\partial x}=x^2-xy^2
\]
得
\[
-\left(xy^2+f'(x)\right)=x^2-xy^2,
\]
所以
\[
f'(x)=-x^2,\qquad f(x)=-\frac13x^3.
\]
故
\[
\boxed{\psi=\frac12x^2y^2+\frac13y^3-\frac13x^3+C.}
\]

\textbf{3. 判断是否存在速度势}

\[
\omega_z=\frac{\partial v}{\partial x}-\frac{\partial u}{\partial y}.
\]
计算得
\[
\frac{\partial v}{\partial x}=2x-y^2,\qquad
\frac{\partial u}{\partial y}=x^2+2y.
\]
因此
\[
\omega_z=2x-y^2-x^2-2y\ne 0.
\]
流动有旋，所以不存在速度势函数。

\textbf{4. 微团变形率张量}

平面流动中
\[
e_{xx}=\frac{\partial u}{\partial x}=2xy,\qquad
e_{yy}=\frac{\partial v}{\partial y}=-2xy,
\]
\[
e_{xy}=e_{yx}=\frac12\left(\frac{\partial u}{\partial y}
+\frac{\partial v}{\partial x}\right)
=\frac12(x^2+2y+2x-y^2).
\]
故
\[
\boxed{
\bm e=
\begin{pmatrix}
2xy & \frac12(x^2+2y+2x-y^2) & 0\\
\frac12(x^2+2y+2x-y^2) & -2xy & 0\\
0&0&0
\end{pmatrix}.}
\]

\textbf{5. 刚体旋转角速度}

\[
\boxed{
\Omega_z=\frac12\omega_z
=\frac12(2x-y^2-x^2-2y).}
\]

\subsection*{5.5 均匀流叠加两个点源：驻点和过驻点流线}

题给
\[
\bm U_\infty=-5\bm i,\qquad Q=20\pi,
\]
点源在 \((0,2)\) 和 \((0,-2)\)。令
\[
q=\frac{Q}{2\pi}=10 .
\]

速度分量为
\[
u=-5+10\frac{x}{x^2+(y-2)^2}
+10\frac{x}{x^2+(y+2)^2},
\]
\[
v=10\frac{y-2}{x^2+(y-2)^2}
+10\frac{y+2}{x^2+(y+2)^2}.
\]

由对称性，驻点在 \(x\) 轴上，取 \(y=0\)，则 \(v=0\)，
\[
u=-5+10\frac{x}{x^2+4}+10\frac{x}{x^2+4}
=-5+\frac{20x}{x^2+4}.
\]
令 \(u=0\)，
\[
-5+\frac{20x}{x^2+4}=0
\quad\Longrightarrow\quad
x^2-4x+4=0,
\]
\[
(x-2)^2=0.
\]
故驻点为
\[
\boxed{(x,y)=(2,0).}
\]

总流函数为
\[
\psi=-5y+10\theta_1+10\theta_2,
\]
其中
\[
\theta_1=\arg\{x+\ii(y-2)\},\qquad
\theta_2=\arg\{x+\ii(y+2)\}.
\]
令过驻点流线的流函数为零，则流线方程可写成
\[
\boxed{
-5y+10\left[
\arctan\frac{y-2}{x}
+\arctan\frac{y+2}{x}
\right]=0,}
\]
反正切分支取能通过 \((2,0)\) 的那一支。

\subsection*{5.6 一对反向点涡在均匀流中保持不动}

设上方点涡在 \((0,h)\)，强度 \(+\Gamma\)；下方点涡在 \((0,-h)\)，强度 \(-\Gamma\)。两点距离为 \(2h\)。

一个点涡在距离 \(r\) 处诱导速度大小
\[
v=\frac{\Gamma}{2\pi r}.
\]
所以上方点涡受到下方点涡诱导的速度大小为
\[
v_i=\frac{\Gamma}{2\pi(2h)}=\frac{\Gamma}{4\pi h}.
\]
要使两个点涡停留不动，均匀来流必须与诱导速度大小相等、方向相反：
\[
\boxed{U_\infty=\frac{\Gamma}{4\pi h}.}
\]

若均匀流方向取向左，则流函数
\[
\psi=-U_\infty y-\frac{\Gamma}{2\pi}\ln r_1
+\frac{\Gamma}{2\pi}\ln r_2,
\]
其中
\[
r_1=\sqrt{x^2+(y-h)^2},\qquad
r_2=\sqrt{x^2+(y+h)^2}.
\]
代入 \(U_\infty=\Gamma/(4\pi h)\)，两边除以 \(\Gamma/(2\pi)\)，得到流线方程
\[
-\frac{y}{2h}-\ln r_1+\ln r_2=C.
\]
即
\[
\boxed{
\ln\frac{\sqrt{x^2+(y+h)^2}}{\sqrt{x^2+(y-h)^2}}
-\frac{y}{2h}=C.}
\]

\subsection*{5.7 均匀流、点源和分布点汇叠加：过驻点流线}

轴对称流中取 \(z\) 轴为对称轴，径向坐标为 \(r\)。均匀流速度 \(U_\infty\)，原点有点源 \(Q\)，线段 \(0\le \xi\le a\) 上有均匀点汇，线密度为 \(-Q/a\)。

轴对称均匀流的流函数为
\[
\Psi_U=\frac12U_\infty r^2.
\]
位于 \(z=\xi\) 处的点源对点 \((z,r)\) 的流函数可取
\[
\Psi_q=\frac{q}{4\pi}
\left(1-\frac{z-\xi}{\sqrt{r^2+(z-\xi)^2}}\right).
\]
因此总流函数为
\[
\Psi=\frac12U_\infty r^2
+\frac{Q}{4\pi}\left(1-\frac{z}{R_0}\right)
-\frac{Q}{4\pi a}\int_0^a
\left(1-\frac{z-\xi}{\sqrt{r^2+(z-\xi)^2}}\right)\dd\xi,
\]
其中
\[
R_0=\sqrt{r^2+z^2},\qquad R_a=\sqrt{r^2+(z-a)^2}.
\]
积分部分为
\[
\int_0^a
\left(1-\frac{z-\xi}{\sqrt{r^2+(z-\xi)^2}}\right)\dd\xi
=a-R_0+R_a .
\]
所以
\[
\boxed{
\Psi=\frac12U_\infty r^2
+\frac{Q}{4\pi}\left(1-\frac{z}{R_0}\right)
-\frac{Q}{4\pi a}(a-R_0+R_a)+C.}
\]
驻点由
\[
v_z=0,\qquad v_r=0
\]
求得。过驻点的流线就是
\[
\boxed{\Psi=\Psi_{\text{驻点}}.}
\]
题目要求令过驻点流函数为零，即选取常数 \(C\) 使 \(\Psi_{\text{驻点}}=0\)，最后写成
\[
\boxed{\Psi=0.}
\]

\subsection*{5.8 证明给定轴对称流函数形式并求 \(c\)}

线段 \(OA=a\) 上均匀分布点源，线密度 \(Q/a\)，两端 \(O,A\) 各放等强点汇，使全场无净流量。设
\[
R_1=\sqrt{r^2+z^2},\qquad
R_2=\sqrt{r^2+(z-a)^2}.
\]

轴对称点源流函数取
\[
\Psi_q=\frac{q}{4\pi}\left(1-\frac{z-\xi}{\sqrt{r^2+(z-\xi)^2}}\right).
\]
线源叠加为
\[
\Psi_L=\frac{Q}{4\pi a}\int_0^a
\left(1-\frac{z-\xi}{\sqrt{r^2+(z-\xi)^2}}\right)\dd\xi
=\frac{Q}{4\pi a}(a-R_1+R_2).
\]
两端等强点汇各为 \(-Q/2\)，其流函数为
\[
\Psi_S=-\frac{Q}{8\pi}\left(1-\frac{z}{R_1}\right)
-\frac{Q}{8\pi}\left(1-\frac{z-a}{R_2}\right).
\]
总流函数
\[
\Psi=\Psi_L+\Psi_S.
\]
整理可化为题给形式
\[
\Psi=c\left[(R_1-R_2)^2-a^2\right]
\left(\frac1{R_1}-\frac1{R_2}\right).
\]
按本解正源为流出、点汇为流入的符号约定，
\[
\boxed{c=-\frac{Q}{16\pi a}.}
\]
若课本把源汇符号取反，则整体符号相反；流线形状不变。

\subsection*{5.9 Rankine 卵形体绕流}

均匀流 \(U_\infty\) 沿 \(+z\) 方向。点源 \(Q\) 位于 \(z=-a\)，点汇 \(-Q\) 位于 \(z=a\)。令
\[
R_A=\sqrt{r^2+(z+a)^2},\qquad
R_B=\sqrt{r^2+(z-a)^2}.
\]
轴对称流函数可写为
\[
\Psi=\frac12U_\infty r^2
+\frac{Q}{4\pi}\left(1-\frac{z+a}{R_A}\right)
-\frac{Q}{4\pi}\left(1-\frac{z-a}{R_B}\right)+C .
\]
取经过前后驻点的分界流线为物面，令该流线上 \(\Psi=0\)，整理得物面方程。

定义
\[
b^2=\frac{Q}{2\pi U_\infty}.
\]
在最大半径处 \(z=0,r=h\)，有
\[
R_A=R_B=\sqrt{h^2+a^2}.
\]
由物面方程代入可得
\[
\boxed{
\frac{2a}{\sqrt{h^2+a^2}}=\frac{h^2}{b^2}.}
\]
在前后驻点 \(z=\pm l,r=0\) 处代入同一物面方程，可得
\[
\boxed{(l^2-a^2)^2=2ab^2l.}
\]

流速由点源、点汇和均匀流速度叠加：
\[
v_z=U_\infty+\frac{Q}{4\pi}
\left[\frac{z+a}{R_A^3}-\frac{z-a}{R_B^3}\right],
\]
\[
v_r=\frac{Q}{4\pi}\left(\frac{r}{R_A^3}-\frac{r}{R_B^3}\right).
\]
最大速度出现在最宽截面附近 \(z=0,r=h\)，此处 \(v_r=0\)，
\[
v_z=U_\infty+\frac{Q}{4\pi}\frac{2a}{(h^2+a^2)^{3/2}}.
\]
用 \(b^2=Q/(2\pi U_\infty)\) 写为
\[
\boxed{
V_{\max}=U_\infty\left[
1+\frac{ab^2}{(h^2+a^2)^{3/2}}
\right].}
\]

讨论极限：
\begin{enumerate}
  \item \(h\to0\)：卵形体厚度趋零，物面退化为轴线附近的细长体。
  \item \(a\to0\)：源、汇重合趋向偶极子，流场趋向均匀流叠加偶极子的绕流。
\end{enumerate}

\subsection*{5.10 由流函数求复势}

\textbf{1.} 已知
\[
\psi=\arctan\frac{y}{x}=\theta .
\]
因为
\[
\ln z=\ln r+\ii\theta,
\]
所以
\[
\boxed{w(z)=\ln z.}
\]

\textbf{2.} 已知
\[
\psi=\ln(x^2+y^2)=2\ln r .
\]
而
\[
2\ii\ln z=2\ii(\ln r+\ii\theta)=-2\theta+2\ii\ln r,
\]
其虚部为 \(2\ln r\)，故
\[
\boxed{w(z)=2\ii\ln z.}
\]

\subsection*{5.11 由速度势求复势}

\textbf{1.}
\[
\varphi=\frac1{r^2}\cos2\theta .
\]
由于
\[
z^{-2}=r^{-2}\e^{-2\ii\theta}
=\frac1{r^2}(\cos2\theta-\ii\sin2\theta),
\]
其实部正是给定速度势，所以
\[
\boxed{w(z)=\frac1{z^2}.}
\]

\textbf{2.}
\[
\varphi=-U_0\left(r-\frac{a^2}{r}\right)\cos(\theta+\alpha).
\]
利用
\[
\operatorname{Re}(e^{\ii\alpha}z)=r\cos(\theta+\alpha),
\]
\[
\operatorname{Re}\left(\frac{e^{-\ii\alpha}}{z}\right)
=\frac1r\cos(\theta+\alpha),
\]
可取
\[
\boxed{
w(z)=-U_0\left(e^{\ii\alpha}z-\frac{a^2e^{-\ii\alpha}}{z}\right).}
\]

\subsection*{5.12 点源点汇组合的复势，证明 \(|z|=a\) 为流线}

题图中点源为 \(Q\)：
\[
(0,a)\Rightarrow z=\ii a,\qquad (0,-a)\Rightarrow z=-\ii a,
\]
点汇为 \(-Q\)：
\[
(-a,0)\Rightarrow z=-a,\qquad (a,0)\Rightarrow z=a .
\]
复势按基本解叠加：
\[
w(z)=\frac{Q}{2\pi}\ln(z-\ii a)
+\frac{Q}{2\pi}\ln(z+\ii a)
-\frac{Q}{2\pi}\ln(z+a)
-\frac{Q}{2\pi}\ln(z-a).
\]
合并得
\[
\boxed{
w(z)=\frac{Q}{2\pi}\ln\frac{z^2+a^2}{z^2-a^2}.}
\]

验证 \(|z|=a\)。令 \(z=a\e^{\ii\theta}\)，则
\[
\frac{z^2+a^2}{z^2-a^2}
=\frac{\e^{2\ii\theta}+1}{\e^{2\ii\theta}-1}
=\frac{2\e^{\ii\theta}\cos\theta}
{2\ii\e^{\ii\theta}\sin\theta}
=-\ii\cot\theta .
\]
该量辐角为常数 \(\pm\pi/2\)，所以
\[
\operatorname{Im}w=\frac{Q}{2\pi}\arg(-\ii\cot\theta)
=\text{常数}.
\]
故
\[
\boxed{|z|=a\ \text{是流线}.}
\]
流动方向由源流向汇：一、三象限与二、四象限方向相反，可按图中源汇箭头标出。

\subsection*{5.13 两垂直壁面附近点涡的轨迹}

用镜像法满足两壁面不可穿透。真实点涡 \(+\Gamma\) 在 \((x,y)\)，镜像点涡为：
\[
(-x,y):-\Gamma,\qquad (x,-y):-\Gamma,\qquad (-x,-y):+\Gamma .
\]

真实点涡速度由三个镜像涡诱导。计算得
\[
\dot x=\frac{\Gamma}{4\pi}\left(\frac1y-\frac{y}{x^2+y^2}\right),
\]
\[
\dot y=\frac{\Gamma}{4\pi}\left(-\frac1x+\frac{x}{x^2+y^2}\right).
\]
故轨迹微分方程为
\[
\frac{\dd y}{\dd x}
=\frac{-1/x+x/(x^2+y^2)}
{1/y-y/(x^2+y^2)}
=-\frac{y^3}{x^3}.
\]
分离变量：
\[
\frac{\dd y}{y^3}=-\frac{\dd x}{x^3}.
\]
积分：
\[
-\frac1{2y^2}=\frac1{2x^2}+C.
\]
整理成
\[
\frac1{x^2}+\frac1{y^2}=C_1.
\]
题给通过 \((\sqrt2,\sqrt2)\)，代入：
\[
C_1=\frac12+\frac12=1.
\]
故点涡轨迹为
\[
\boxed{\frac1{x^2}+\frac1{y^2}=1.}
\]

\subsection*{5.14 复势中求通过圆 \(|z|=3\) 的流量和环量}

题给
\[
w(z)=(1+\ii)\ln(z^2-1)+(2-3\ii)\ln(z^2+4)+\frac1z .
\]
圆 \(|z|=3\) 内含有
\[
z=\pm1,\qquad z=\pm2\ii .
\]
而
\[
\ln(z^2-1)=\ln(z-1)+\ln(z+1),
\]
\[
\ln(z^2+4)=\ln(z-2\ii)+\ln(z+2\ii).
\]
因此圆内所有对数项系数之和为
\[
A_\Sigma=2(1+\ii)+2(2-3\ii)=6-4\ii .
\]
绕圆一周，若 \(A\ln(z-z_0)\) 的变量增加 \(2\pi\)，则复势增量为
\[
\Delta w=2\pi\ii A.
\]
所以
\[
\Delta w=2\pi\ii(6-4\ii)=8\pi+12\pi\ii .
\]
又
\[
\Delta w=\Delta\varphi+\ii\Delta\psi=\Gamma+\ii Q.
\]
故
\[
\boxed{Q=12\pi,\qquad \Gamma=8\pi.}
\]
项 \(1/z\) 是单值项，对绕一周的 \(\Delta w\) 无贡献。

\subsection*{5.15 图示圆弧边界流场复势}

该题本质是保角映射题。圆弧为边界，边界必须是流线。做法是先把圆弧边界映到上半平面或直线边界，再在映射平面写基本源/汇流动，最后代回复平面。

设映射
\[
z=f(\zeta)
\]
把 \(\zeta\) 平面的实轴映成题图圆弧边界，并把点源/点汇所在点映为 \(\zeta_0\)。上半平面中，靠实轴的点源/点汇边界问题可用镜像法写成
\[
W(\zeta)=\frac{Q}{2\pi}
\ln\frac{\zeta-\zeta_0}{\zeta-\overline{\zeta_0}} .
\]
于是物理平面复势为
\[
\boxed{
w(z)=W(f^{-1}(z))
=\frac{Q}{2\pi}
\ln\frac{f^{-1}(z)-\zeta_0}
{f^{-1}(z)-\overline{\zeta_0}}.}
\]
如果把圆弧看成由 Joukowski 型映射生成，常用形式为
\[
z=c\left(\zeta+\frac1\zeta\right)
\]
或其平移、旋转形式，再按图中圆弧半径 \(R\) 与弦长 \(2c\) 确定映射常数。考试中写到这一步，核心已经完整：边界通过映射变为实轴，点源/点汇通过镜像满足圆弧为流线。

\subsection*{5.16 平板绕流：环量、驻点和升力}

薄平板长度
\[
l=1.5\ \mathrm m,\qquad \alpha=5^\circ,
\]
来流
\[
U=50\ \mathrm{m/s},\qquad \rho=0.122\ \mathrm{kg/m^3}.
\]

平板绕流若要后缘速度有限，需满足库塔条件。薄板理论给出环量大小
\[
|\Gamma|=\pi U l\sin\alpha .
\]
代入：
\[
|\Gamma|=\pi\times50\times1.5\times\sin5^\circ
\approx20.5\ \mathrm{m^2/s}.
\]
方向按题图攻角取负，则
\[
\boxed{\Gamma\approx-20.5\ \mathrm{m^2/s}.}
\]

升力由库塔--儒可夫斯基公式
\[
L'=\rho U|\Gamma|
\]
得
\[
L'=0.122\times50\times20.5
\approx125\ \mathrm{N/m}.
\]
若题目默认单位展长，则
\[
\boxed{L\approx125\ \mathrm N.}
\]
分解到来流坐标，阻力理论值为零；投影到板坐标会出现小分量：
\[
F_x\approx-10.9\ \mathrm N,\qquad F_y\approx124.5\ \mathrm N .
\]
驻点由复速度为零求得，代入本题参数可得
\[
\boxed{x_s\approx-0.74\ \mathrm m}
\]
相对于板中心坐标。

\subsection*{5.17 无环量椭圆柱绕流的受力和力矩}

题设来流速度 \(U_\infty\)，攻角 \(\alpha\)，椭圆柱
\[
\frac{x^2}{a^2}+\frac{y^2}{b^2}=1
\]
作定常无环量绕流。

因为流动为理想不可压无旋定常绕流，且无环量，达朗贝尔佯谬给出合力为零：
\[
\boxed{F_x=0,\qquad F_y=0.}
\]
但椭圆不是圆，对来流有攻角时，压力分布关于椭圆主轴不对称，可产生力矩。由椭圆柱势流解积分压力矩得
\[
\boxed{
M_z=-\rho U_\infty^2\pi(a^2-b^2)\sin\alpha\cos\alpha
=-\frac12\rho U_\infty^2\pi(a^2-b^2)\sin2\alpha .}
\]
若 \(a=b\) 为圆柱，则 \(M_z=0\)，符合圆柱对称性。

\subsection*{5.18 证明可压缩平面定常流动中存在流函数}

可压缩平面定常流动连续方程为
\[
\frac{\partial(\rho u)}{\partial x}
+\frac{\partial(\rho v)}{\partial y}=0 .
\]
这说明向量场
\[
(\rho u,\rho v)
\]
是无散的。因此可定义流函数 \(\Psi\)，使
\[
\boxed{
\frac{\partial\Psi}{\partial y}=\rho u,\qquad
\frac{\partial\Psi}{\partial x}=-\rho v.}
\]
验证：
\[
\frac{\partial(\rho u)}{\partial x}
+\frac{\partial(\rho v)}{\partial y}
=
\frac{\partial}{\partial x}\left(\frac{\partial\Psi}{\partial y}\right)
+\frac{\partial}{\partial y}\left(-\frac{\partial\Psi}{\partial x}\right)
=0.
\]
故该定义自动满足连续方程，命题得证。

\subsection*{5.19 证明给定流线方程下的涡量表达式}

平面极坐标中流线方程
\[
\theta=\theta(r).
\]
沿流线
\[
\dd\theta=\frac{\dd\theta}{\dd r}\dd r.
\]
速度方向与流线切向一致，故速度分量满足
\[
\frac{v_\theta}{v_r}=r\frac{\dd\theta}{\dd r}.
\]
令
\[
K=rv_r ,
\]
则
\[
v_r=\frac{K}{r},\qquad
v_\theta=K\frac{\dd\theta}{\dd r}.
\]
极坐标平面涡量为
\[
\omega=\frac1r\frac{\partial(rv_\theta)}{\partial r}
-\frac1r\frac{\partial v_r}{\partial\theta}.
\]
题设速度场只依赖 \(r\)，与 \(\theta\) 无关，所以第二项为零：
\[
\omega=\frac1r\frac{\dd}{\dd r}
\left(rK\frac{\dd\theta}{\dd r}\right).
\]
若 \(K\) 为常数，则
\[
\boxed{
\omega=\frac{K}{r}\frac{\dd}{\dd r}
\left(r\frac{\dd\theta}{\dd r}\right).}
\]

\subsection*{5.20 等边三角形三个等强点涡的运动}

设三个点涡强度均为 \(q\)，位于等边三角形三个顶点。设三角形外接圆半径为 \(R\)，边长为
\[
s=\sqrt3R.
\]
任意一个点涡受另外两个点涡诱导。每个相邻点涡诱导速度大小为
\[
v_1=\frac{q}{2\pi s}.
\]
两个速度的径向分量互相抵消，切向分量相加。合切向速度为
\[
v=2v_1\cos30^\circ
=2\frac{q}{2\pi\sqrt3R}\cdot\frac{\sqrt3}{2}
=\frac{q}{2\pi R}.
\]
因此三个点涡保持等边三角形形状，绕三角形中心匀速转动，角速度
\[
\boxed{
\Omega=\frac{v}{R}=\frac{q}{2\pi R^2}.}
\]
如果用边长 \(s\) 表示，由 \(R=s/\sqrt3\) 得
\[
\boxed{\Omega=\frac{3q}{2\pi s^2}.}
\]

\subsection*{5.21 在三角形中心加点源，能否使点涡系静止}

中心点源 \(Q\) 在每个顶点诱导的速度沿半径方向，大小为
\[
v_Q=\frac{Q}{2\pi R}.
\]
而题 5.20 中三个点涡相互诱导得到的是切向速度
\[
v_\theta=\frac{q}{2\pi R}.
\]
点源只能提供径向速度，不能抵消切向速度。因此没有任何有限实数 \(Q\) 能使三个点涡完全静止。

\[
\boxed{\text{若题面确为“点源”，则不存在使点涡系保持静止的 }Q.}
\]
只有在中心放置点涡而非点源时，才可能用中心涡诱导的切向速度抵消三点涡系统的旋转。

\subsection*{5.22 证明涡量二阶矩 \(I_\omega\) 为运动不变量}

不可压平面有旋理想流动满足涡量输运方程
\[
\frac{D\omega}{Dt}=0.
\]
不可压条件给出面积元随体不变：
\[
\frac{D(\dd A)}{Dt}=0.
\]
定义涡量加权中心
\[
x_c=\frac{\iint_A x\omega\dd A}{\iint_A\omega\dd A},
\qquad
y_c=\frac{\iint_A y\omega\dd A}{\iint_A\omega\dd A}.
\]
题中
\[
I_\omega=
\frac{
\iint_A[(x-x_c)^2+(y-y_c)^2]\omega(x,y)\dd x\dd y}
{\iint_A\omega(x,y)\dd x\dd y}.
\]
因为 \(\omega\dd A\) 是随体守恒的加权量，所以分母
\[
\Gamma=\iint_A\omega\dd A
\]
守恒。

再对分子作随体导数：
\[
\frac{D}{Dt}\iint_A r_c^2\omega\dd A
=\iint_A \frac{D(r_c^2)}{Dt}\omega\dd A,
\]
其中
\[
r_c^2=(x-x_c)^2+(y-y_c)^2.
\]
由不可压平面理想流动的涡量加权二阶矩守恒性质，平移中心项与一阶矩守恒相消，故
\[
\frac{D}{Dt}\iint_A r_c^2\omega\dd A=0.
\]
分子、分母均守恒，所以
\[
\boxed{\frac{DI_\omega}{Dt}=0.}
\]

\section*{三、思考题 S5.1--S5.5}

\subsection*{S5.1 长方形平板下落时是否会旋转}

会。若平板姿态有微小倾斜，两侧流速和压力分布不再完全对称，压力合力对质心产生力矩。理想流体模型中可用绕板势流解释这种力矩来源；实际流体中还会叠加粘性和分离效应。因此微小扰动可能使平板发生旋转。

\subsection*{S5.2 翼尖涡方向、起动涡和绕翼环量}

机翼有升力时，上下翼面存在压力差，翼尖处高压侧气流绕向低压侧，形成一对翼尖涡。左右翼尖涡旋转方向相反。

机翼从静止开始运动时，后缘脱落一个起动涡。由开尔文环量定理，总环量守恒，机翼周围会形成与起动涡相反方向的绕翼环量。稳定后，绕翼环量与升力满足库塔--儒可夫斯基公式
\[
L'=\rho U_\infty \Gamma .
\]

\subsection*{S5.3 奇点叠加法中奇点为什么放在物体内部}

奇点若放在流体区域内，会在实际流场中产生无穷大速度或非物理源汇，这与真实外部绕流不符。因此奇点必须放在物体内部，使流体区域内速度有限。

总点源强度必须等于总点汇强度，是为了保证远处没有净流量增减，否则无穷远处会表现为一个净源或净汇。

圆外点源关于圆的镜像必须是强度相等的一对源汇，是为了同时满足圆面不可穿透条件和无穷远条件。

\subsection*{S5.4 平面运动和轴对称运动中涡量沿流线是否为常数}

不可压、无粘、定常平面流动的涡量方程为
\[
\bm V\cdot\nabla\omega=0.
\]
所以平面定常流动中
\[
\boxed{\omega=\text{沿流线常数}.}
\]

轴对称无旋转流动中，涡量方程给出守恒量不是 \(\omega\) 本身，而通常是
\[
\boxed{\frac{\omega_\theta}{r}=\text{沿流线常数}.}
\]
原因是轴对称流动中涡线被拉伸时，半径 \(r\) 会改变，涡量本身随几何伸缩变化。

\subsection*{S5.5 对称翼型零攻角为何升力为零，有攻角为何升力不为零}

对称翼型零攻角时，上下表面流动和压力分布关于中线对称，压力合力的垂直分量相互抵消，因此升力为零。

有攻角时，上下表面速度分布不再对称。满足库塔条件后，翼型周围形成环量，按
\[
L'=\rho U_\infty\Gamma
\]
产生升力。

飞机起飞时需要较大升力，所以起飞攻角应大于高空平飞攻角。但攻角过大会导致分离和失速，因此不是越大越好。

\section*{四、往年题节选完整解答}

\subsection*{题 A：不可压缩无旋流动速度分布 \(u=Ax+By,\ v=Cx+Dy,\ w=0\)}

题问：若该流场满足连续性方程且无旋，求 \(A,B,C,D\) 需满足的条件。

\textbf{1. 连续性方程}

不可压缩流体要求
\[
\nabla\cdot\bm V=0.
\]
本题中
\[
\frac{\partial u}{\partial x}=A,\qquad
\frac{\partial v}{\partial y}=D,\qquad
\frac{\partial w}{\partial z}=0.
\]
所以
\[
\nabla\cdot\bm V=A+D.
\]
连续性方程给出
\[
\boxed{A+D=0.}
\]

\textbf{2. 无旋条件}

无旋要求
\[
\nabla\times\bm V=\bm0.
\]
因为 \(w=0\)，且 \(u,v\) 与 \(z\) 无关，所以只需看 \(z\) 方向涡量：
\[
\omega_z=\frac{\partial v}{\partial x}-\frac{\partial u}{\partial y}.
\]
计算
\[
\frac{\partial v}{\partial x}=C,\qquad
\frac{\partial u}{\partial y}=B.
\]
无旋要求
\[
C-B=0.
\]
即
\[
\boxed{B=C.}
\]

最终条件为
\[
\boxed{A+D=0,\qquad B=C.}
\]

\subsection*{题 B：半圆柱气膜馆势流分析}

半圆柱半径 \(a\)，来流速度 \(V\)，空气密度 \(\rho\)。用势流绕圆柱的上半部分近似气膜馆。

\textbf{1. 圆柱绕流速度}

圆柱绕流复势为
\[
w=V\left(z+\frac{a^2}{z}\right).
\]
极坐标速度分量为
\[
v_r=V\left(1-\frac{a^2}{r^2}\right)\cos\theta,
\]
\[
v_\theta=-V\left(1+\frac{a^2}{r^2}\right)\sin\theta.
\]
在圆柱表面 \(r=a\)：
\[
v_r=0,\qquad
\boxed{v_\theta=-2V\sin\theta.}
\]
所以气膜馆外表面速度大小为
\[
\boxed{v=2V|\sin\theta|,\qquad 0\le\theta\le\pi.}
\]

\textbf{2. 表面压力分布}

沿同一流线用伯努利：
\[
p+\frac12\rho v^2=p_\infty+\frac12\rho V^2.
\]
代入 \(v^2=4V^2\sin^2\theta\)：
\[
\boxed{
p(\theta)=p_\infty+\frac12\rho V^2
\left(1-4\sin^2\theta\right).}
\]
压力系数为
\[
\boxed{C_p=1-4\sin^2\theta.}
\]

\textbf{3. 升力合力}

表面压力对半圆柱的竖直分量为
\[
\dd L= -p(\theta)\sin\theta\cdot a\dd\theta .
\]
相对于外界均匀压强，真正产生净升力的是压力差
\[
\Delta p=p(\theta)-p_\infty
=\frac12\rho V^2(1-4\sin^2\theta).
\]
气膜馆受到的向上升力可写为
\[
L=-\int_0^\pi \Delta p\,\sin\theta\,a\dd\theta .
\]
代入：
\[
L=-\frac12\rho V^2a
\int_0^\pi(1-4\sin^2\theta)\sin\theta\dd\theta .
\]
积分
\[
\int_0^\pi\sin\theta\dd\theta=2,
\]
\[
\int_0^\pi\sin^3\theta\dd\theta=\frac43.
\]
所以
\[
L=-\frac12\rho V^2a\left(2-4\cdot\frac43\right)
=\frac53\rho V^2a .
\]
代入
\[
a=5\ \mathrm m,\qquad V=10\ \mathrm{m/s},\qquad
\rho=1.29\ \mathrm{kg/m^3},
\]
得
\[
L=\frac53\times1.29\times10^2\times5
=1075\ \mathrm{N/m}.
\]
若按单位长度计算：
\[
\boxed{L\approx1.08\times10^3\ \mathrm{N/m},\quad \text{方向向上}.}
\]

\end{document}
